\section{The Physical World}
\label{sec:physical}

The physical world is the most familiar lawful region of the substrate, the part settled into the
shared, stable pattern we all inhabit. The program has been to check, in a finite testbed, that the major features of
physics actually emerge from the one rule on the mesh, rather than to assume them. What follows
is the picture in plain terms, with the honest standing of each part noted. Throughout, the
move is the same: a feature of our world is a large-scale pattern in the web.

\paragraph{Space and time.}
Space is the shape of the web seen from far away. Up close there are only vibes and the notes
between them, and distance is how many notes one must cross. Stepped back, the pattern of nearness
behaves like a smooth geometry with a definite number of dimensions, recovered from the
relations alone, and the recovered dimension sharpens as the web grows. Time is the depth of
the web's before-and-after, the length of the longest chain of ``and then.'' The flow of time
and energy is a local operator on the grown web that is, all at once, local, bounded below so
the world does not fall forever, and finite-speed, giving a light cone. The arrow of time is
simply that the web only grows.

\paragraph{Matter and particles.}
A particle is not a tiny ball but a persistent pattern, a knot in the web that holds its shape
as the universe evolves around it, the way a vortex keeps its form in a flowing fluid. Different
particles are different stable knots. Spin is a feature of the knot that cannot be smoothly
undone, locked in and quantized. Mass is a gap, the smallest energy needed to excite the
pattern, and a massive pattern obeys the relativistic relation between energy and momentum. The
field types sorted by spin, the scalar, the matter spinor, the light-carrying vector, and the
tensor of gravity, are all present as patterns, and mass is generated, not inserted, by a
Higgs-like mechanism in which a background field settles to a nonzero value.

\paragraph{Forces.}
A force is the twist a tone picks up as it crosses a note, the structure physics already uses
for electromagnetism and the nuclear forces. In the testbed this reproduces confinement, the
reason quarks are never seen alone, along with the bookkeeping relation between the topology of
a force field and the particle content, and a condensate. The operators for light and for
gravity are not put in by hand but appear as the gentle-vibration limit of the one action.

\paragraph{Gravity.}
Gravity is the web's own curvature and how it responds to what sits in it. In the everyday,
slow, weak limit the model returns Newton's law. Beyond that, the linearized Einstein equation
appears as a genuine law of motion that conserves energy and momentum, reduces to Newton when
things are slow, and carries gravitational waves at the speed of light, and the corresponding
graviton built directly on the discrete lattice is gauge-invariant, massless, and spin-two. The
full nonlinear equation holds in its cosmological form, where its three parts stay consistent
through the Bianchi identity. Because the web is discrete there is a smallest length, so
curvature cannot become infinite, and the singularities of standard gravity are softened into
something finite. Black-hole entropy comes out proportional to horizon area rather than volume,
because entanglement across a boundary scales with the boundary, and the horizon carries a
temperature, a thermal spectrum at the Hawking value, that rises as the hole shrinks.

\paragraph{The quantum.}
This is where the model is most interesting and where it is most careful. A quantum amplitude is
a smooth complex number, while a vibe is three-valued, so the amplitude cannot live on a vibe.
It is an average over a patch of many vibes: its size a density of how many are excited
together, its phase the patch's shared rhythm. With that in view the pillars appear. Evolution
conserves probability and a disturbance spreads in the wave-like way that signals coherence.
Interference is present, possibilities adding and canceling. The Bell correlations are
reproduced, and the model says why: because there is no hidden state, the measuring device and
the measured thing are made of the same web and share a common past, so they are not
independent, and that shared substrate carries the correlation with no signal traveling faster
than light. The Born rule, that probability is the amplitude squared, is derived rather than
postulated, because the amplitude's size is the square root of a density of vibes, so fair
sampling of the substrate gives the squared amplitude, with the exponent two singled out.

\paragraph{Cosmology and the dark sector.}
A universe that grows forever naturally expands, and the testbed shows expansion from the local
growth rule, with a brief early burst and a graceful settling. Dark energy comes out near the
right, famously tiny, magnitude as the natural fluctuation of the cosmological constant, tied to
the fluctuating spacetime volume, which at the size of our universe lands close to the observed
value. Dark matter has a candidate mechanism, a small long-range correction to gravity that
flattens galaxy rotation curves without a new particle.

\paragraph{What this amounts to.}
Taken together, the substrate reproduces a wide span of known physics from one rule. We are
explicit that this is a finite simulation and a model, not a derivation of nature. Some pieces
are solid, some are first rungs on deep problems, and the deepest quantum and gravitational
questions are partly open. The point is not that everything is settled, but that the physical
world, in its main features, can be made to come out of an experiential base, which is what a
strict monism must show.
